% \section{Declarations}
\section{声明}

% \paragraph{Availability of Data and Material.}
\paragraph{数据及材料的可用性。}
% Our model and the involved dataset will be publicly available. Meanwhile, we will also release the instruction construction pipeline. The used fMRI signal and instruction pairs will be released for the research community to further explore this task.
我们的模型和相关数据集将公开发布。同时，我们也将开源指令构建的流水线代码。此外，我们还将发布所使用的 fMRI 信号与指令对，以供研究社区进一步探索这一任务。

% \paragraph{Competing Interests.}
\paragraph{竞争利益。}
% All authors certify that they have no affiliation or involvement in any organization or entity with any financial or non-financial interest in the subject matter or materials discussed in this manuscript.
所有作者均在此声明，未隶属于任何组织或实体，也未参与任何涉及本文所讨论的主题或材料且存在财务或非财务利益关系的组织或实体。

% \paragraph{Authors' Contributions.}
\paragraph{作者贡献。}
% Yasheng Sun coordinated the teamwork and proposed the overall method. Bohan Li trained the fMRI-to-Image module and conducted comparison results.
孙亚胜（Yasheng Sun）协调了团队工作并提出了整体方法。李博涵（Bohan Li）训练了 fMRI 到图像的生成模块并主导了对比实验结果的整理。
% Mingchen Zhuge initialized the idea of multi-modal learning with human brain activity signals and drafted the current version of the abstract and introduction. Prof. Fan manage the whole project. Prof. Fan, Prof. Salaman, Prof. Fahad, and Prof. Koike had insightful discussions and provided help polishing the manuscript.
诸葛鸣晨（Mingchen Zhuge）提出了结合人类大脑活动信号进行多模态学习的初步构想，并起草了当前版本的摘要与引言。范教授（Prof. Fan）管理了整个项目。范教授、Salaman 教授、Fahad 教授以及小池（Koike）教授参与了深刻的讨论，并在论文的润色方面提供了帮助。
% All authors read and approved the final manuscript.
所有作者均已阅读并同意最终定稿。

% \paragraph{Acknowledgements.}
\paragraph{致谢。}
% The authors express their gratitude to the anonymous reviewers and the editor, whose valuable feedback greatly improved the quality of this manuscript.
作者在此对匿名审稿人和编辑表示由衷的感谢，他们宝贵的反馈意见极大地提升了本稿件的质量。